\setcounter{subsection}{4}

\subsection{Eigenvalues and Eigenvectors}

\begin{apex}[5a42]
    Define $T \in \lin{\f^n}$ by $T(x_1, x_2, x_3, \ldots, x_n) = (x_1, 2x_2, 3x_3, \ldots, nx_n)$.
    \begin{subexercises}
        \item Find all eigenvalues and eigenvectors of $T$.
        \item Find all subspaces of $\f^n$ that are invariant under $T$.
    \end{subexercises}
\end{apex}

\begin{csp}
    Let $e_1$, $\dots$, $e_n$ be the standard basis of $\mathbb{F}^n$.
    \begin{subexercises}
        \item For each $k \in \{1, \dots, n\}$, we have $Te_k=ke_k$.
		Hence, the eigenvalues of $T$ are $1$, $2$, $\dots$, $n$ and the eigenvectors corresponding to $k$ are the nonzero scalar multiples of $e_k$.
		Note that this list of eigenvalues (and eigenvectors) is exhaustive because $\dim \left(\mathbb{F}^n \right)=n$.
        \item Let $U \subseteq \mathbb{F}^n$ be an invariant subspace under $T$.
		Take any $u=\sum^{n}_{i=1} \alpha_i e_i \in U$ and suppose $\alpha_k \neq 0$ for some $k$.
		By Lagrange interpolation, there exists a polynomial $p$ such that $p(i)=\delta_{i,k}$ for all $i \in \{1, \dots, n\}$.
		Since $U$ is invariant under $T$, it is also invariant under $p(T)$, and so
        \[p(T)u=\sum^{n}_{i=1} \alpha_i p(T) e_i=\sum^{n}_{i=1} \alpha_i p(i) e_i=\alpha_k e_k \in U \implies e_k \in U.\]
        Repeating this argument for every index with a nonzero coefficient in some element of $U$ shows that $U$ contains each of the basis vectors $e_1$, $\dots$, $e_n$.
		Hence, $U=\operatorname{span} \{e_i: i \in S\}$ for some $S \subseteq \{1, \dots, n\}$.

        Conversely, any such span is invariant because $Te_i=ie_i$ belongs to that span.
		There are $2^n$ such subspaces, and they exhaust all invariant subspaces of $T$. \qh
    \end{subexercises}
\end{csp}